\begin{textsample}{POS Dim 1 – ro – Score 60.00 – ro/ro\_em\_42.txt}  \label{ex:f1_pos_134}
\textbf{Competência} Específica 1 Utilizar estratégias, conceitos e procedimentos matemáticos para interpretar situações em diversos contextos, sejam atividades cotidianas, sejam fatos das Ciências da Natureza e Humanas, das questões socioeconômicas ou tecnológicas, \textbf{divulgados} por diferentes meios, de modo a contribuir para uma formação geral.
A \textbf{competência} 1 apresenta a Matemática como um corpo de conhecimentos a serviço de outras áreas do conhecimento e, por isso, colabora para a formação integral do estudante.
O conhecimento de estratégias, conceitos e procedimentos matemáticos, sempre levando em consideração o contexto em que a situação está inserida, está associado ao domínio da \textbf{competência}.
Ademais, a \textbf{compreensão} do que se deseja determinar, de acordo com cada situação de aprendizagem, \textbf{exige} a combinação de vários conhecimentos de modo apropriado para que seja possível colocar esse conjunto de ideias em ação, monitorando estratégias selecionadas em cada situação e analisando sua eficiência ; bem como a leitura e interpretação de textos verbais, desenhos técnicos, gráficos e imagens.
Percebe -se, portanto, que é uma \textbf{competência} relacionada à preparação dos jovens para construir e realizar Projetos de Vida, destacando -se, ainda, a sua relação com a \textbf{Competência} Geral 2 da BNCC, no que se refere ao exercício da curiosidade intelectual que utiliza o conhecimento para investigar, refletir e criar soluções em diferentes situações.
Habilidades ( EM13MAT101 ).
Interpretar criticamente situações \textbf{econômicas}, sociais e fatos relativos às Ciências da Natureza que envolvam a variação de grandezas, pela \textbf{análise} dos gráficos das funções representadas e das taxas de variação, com ou sem apoio de \textbf{tecnologias} \textbf{digitais}. ( EM13MAT102 ).
Analisar tabelas, gráficos e amostras de pesquisas estatísticas apresentadas em relatórios \textbf{divulgados} por diferentes meios de comunicação, identificando, quando for o caso, inadequações que possam \textbf{induzir} a \textbf{erros} de interpretação, como escalas e amostras não apropriadas. ( EM13MAT103 ).
Interpretar e compreender textos científicos ou \textbf{divulgados} pelas mídias, que empregam unidades de medida de diferentes grandezas e as conversões possíveis entre elas, adotadas ou não pelo Sistema Internacional ( SI ), como as de armazenamento e velocidade de transferência de dados, ligadas aos avanços tecnológicos. ( EM13MAT104 ).
Interpretar taxas e índices de natureza socioeconômica ( índice de desenvolvimento humano, taxas de inflação, entre outros ), investigando os processos de cálculo desses números, para analisar criticamente a realidade e produzir \textbf{argumentos}. ( EM13MAT105 ).
Utilizar as noções de transformações isométricas ( translação, reflexão, rotação e composições destas ) e transformações homotéticas para construir \textbf{figuras} e analisar elementos da natureza e diferentes produções humanas ( fractais, construções civis, obras de arte, entre outras ). ( EM13MAT106 ).
Identificar situações da vida cotidiana nas quais seja necessário fazer escolhas levando -se em conta os riscos probabilísticos ( usar este ou aquele \textbf{método} contraceptivo, optar por um tratamento médico em detrimento de outro etc. ).
Objetos do conhecimento - Funções : interpretação de gráficos e de expressões algébricas.
Sistemas e unidades de medida : leitura e conversão de unidades de grandezas diversas.
Variação de grandezas, como velocidade, concentração, taxas de crescimento ou decrescimento de populações, índices \textbf{econômicos} etc.
Estatística : gráficos ( e \textbf{infográficos} ), medidas de tendência \textbf{central} e de dispersão. - Conceitos estatísticos : população e amostragem.
Gráficos utilizados pela estatística : elementos de um gráfico.
Confiabilidade de fontes de dados.
Correção no traçado de gráficos estatísticos.
Medidas de tendência \textbf{central} e de dispersão. - Funções : representação gráfica e algébrica.
Sistema Internacional de Medidas : principais unidades e conversões.
Bases de sistemas de contagem ( base decimal, base binária, base sexagesimal etc. ).
Principais unidades de armazenamento de dados na informática ( bit, byte, kilobyte, megabyte, gigabyte etc. ) e transferência de dados ( Mbps, Kbps, Gbps etc. ). - Estatística : pesquisa e organização de dados.
Porcentagens : cálculo de índices, taxas e coeficientes.
Estatística : interpretação de gráficos, medidas de tendência \textbf{central} e medidas de dispersão. - Geometria das Transformações : isometrias ( reflexão, translação e rotação ) e homotetias ( ampliação e redução ).
Noções de geometria dos fractais.
Porcentagem : cálculo de taxas, índices e coeficientes.
Probabilidade \textbf{simples} e condicional. \textbf{Eventos} sucessivos, mutuamente exclusivos e não mutuamente exclusivos.
Estatística : distribuição estatística, distribuição normal e medidas de posição ( mediana, quartis, decis e percentis ). \textbf{Competência} Específica 2 Propor ou participar de ações para investigar desafios do mundo \textbf{contemporâneo} e tomar decisões éticas e socialmente responsáveis, com base na \textbf{análise} de \textbf{problemas} sociais, como os voltados a situações de saúde, sustentabilidade, das implicações da \textbf{tecnologia} no mundo do trabalho, entre outros, \textbf{mobilizando} e articulando conceitos, procedimentos e linguagens próprios da Matemática.
A \textbf{competência} 2 propõe a inserção do estudante como personagem atuante em sua comunidade local e no mundo \textbf{globalizado}.
Para \textbf{tanto}, as ações propostas fazem referência à \textbf{capacidade} de ser parte de algo, compartilhar saberes com o outro e colaborar conjuntamente para a produção de algo.
Outrossim, o papel da investigação pressupõe a observação dos desafios presentes na comunidade local/global dos estudantes, a elaboração de hipóteses que os descrevam, o tratamento dos dados associados à situação envolvida, a \textbf{análise} dos resultados obtidos e, por fim, a tomada de decisão a partir das conclusões obtidas.
Portanto, ao desenvolver essa \textbf{competência}, pode -se afirmar que o estudante avança em relação ao entendimento de que os Projetos de Vida não são apenas no âmbito profissional, mas também nas dimensões pessoal e social/cidadã.
Habilidades ( EM13MAT201 ).
Propor ou participar de ações adequadas às \textbf{demandas} da região, preferencialmente para sua comunidade, envolvendo medições e cálculos de perímetro, de área, de volume, de \textbf{capacidade} ou de massa. ( EM13MAT202 ).
Planejar e executar pesquisa amostral sobre questões \textbf{relevantes}, usando dados coletados diretamente ou em diferentes fontes, e comunicar os resultados por meio de relatório contendo gráficos e interpretação das medidas de tendência \textbf{central} e das medidas de dispersão ( amplitude e desvio padrão ), utilizando ou não recursos tecnológicos. ( EM13MAT203 ).
Aplicar conceitos matemáticos no planejamento, na execução e na \textbf{análise} de ações envolvendo a utilização de \textbf{aplicativos} e a criação de planilhas ( para o controle de orçamento familiar, simuladores de cálculos de juros \textbf{simples} e compostos, entre outros ), para tomar decisões.
Objetos do conhecimento - Conceitos e procedimentos de geometria métrica.
Sistema métrico decimal e unidades não convencionais. - Funções, fórmulas e expressões algébricas.
Conceitos \textbf{simples} de Estatística Descritiva.
Medidas de tendência \textbf{central} ( média, moda e mediana ). - Medidas de dispersão ( amplitude, desvio padrão e coeficiente de variância ).
Gráficos estatísticos ( histogramas e polígonos de frequência ). - Distribuição normal.
Cálculos envolvendo porcentagens. - Conceitos de matemática financeira ( juros \textbf{simples}, compostos, taxas de juros etc. ). - Alguns sistemas de amortização e noções de fluxo de caixa.
Funções : exponenciais e logarítmicas. \textbf{Competência} Específica 3 Utilizar estratégias, conceitos, definições e procedimentos matemáticos para interpretar, construir modelos e resolver \textbf{problemas} em diversos contextos, analisando a plausibilidade dos resultados e a adequação das soluções propostas, de modo a construir argumentação consistente.
A \textbf{competência} 3 está relacionada ao chamado “ fazer matemático ”, ou seja, está intimamente ligada à essência da Matemática, cuja ação está pautada na resolução de situações-problema.
Por esse motivo, deixa claro que os conceitos e procedimentos matemáticos somente terão significado caso os estudantes possam utilizá -los para solucionar os desafios com que se deparam.
Com efeito, \textbf{salienta} -se que a referida \textbf{competência} não se \textbf{restringe} apenas à resolução de \textbf{problemas}, mas também trata de sua elaboração.
Isso \textbf{posto}, revela uma concepção mais abrangente da resolução de \textbf{problemas}, além da \textbf{mera} aplicação de um conjunto de regras.
Outro grande destaque refere -se à modelagem matemática como a construção de modelos matemáticos, que sirvam para generalizar ideias ou para descrever situações semelhantes.
Essa \textbf{competência} tem estreita relação com a \textbf{Competência} Geral 2 da BNCC, no sentido da \textbf{capacidade} de formular e resolver \textbf{problemas}, e com a \textbf{Competência} Geral 4, que reforça a importância de saber utilizar as diferentes linguagens para expressar ideias e informações para a comunicação mútua.
Habilidades ( EM13MAT301 ).
Resolver e elaborar \textbf{problemas} do cotidiano, da Matemática e de outras áreas do conhecimento, que envolvem equações lineares simultâneas, usando técnicas algébricas e gráficas, com ou sem apoio de \textbf{tecnologias} \textbf{digitais}. ( EM13MAT302 ).
Construir modelos empregando as funções polinomiais de 1º ou 2º graus, para resolver \textbf{problemas} em contextos diversos, com ou sem apoio de \textbf{tecnologias} \textbf{digitais}. ( EM13MAT303 ).
Interpretar e comparar situações que envolvam juros \textbf{simples} com as que envolvem juros compostos, por meio de representações gráficas ou \textbf{análise} de planilhas, destacando o crescimento linear ou exponencial de cada caso. ( EM13MAT304 ).
Resolver e elaborar \textbf{problemas} com funções exponenciais nos quais seja necessário compreender e interpretar a variação das grandezas envolvidas, em contextos como o da Matemática Financeira, entre outros. ( EM13MAT305 ).
Resolver e elaborar \textbf{problemas} com funções logarítmicas nos quais seja necessário compreender e interpretar a variação das grandezas envolvidas, em contextos como os de abalos sísmicos, pH, radioatividade, Matemática Financeira, entre outros. ( EM13MAT306 ).
Resolver e elaborar \textbf{problemas} em contextos que envolvem fenômenos periódicos reais ( ondas sonoras, fases da lua, movimentos cíclicos, entre outros ) e comparar suas representações com as funções seno e cosseno, no plano cartesiano, com ou sem apoio de \textbf{aplicativos} de álgebra e geometria. ( EM13MAT307 ).
Empregar diferentes \textbf{métodos} para a obtenção da medida da área de uma superfície ( reconfigurações, aproximação por cortes etc. ) e deduzir expressões de cálculo para aplicá -las em situações reais ( como o remanejamento e a distribuição de plantações, entre outros ), com ou sem apoio de \textbf{tecnologias} \textbf{digitais}. ( EM13MAT308 ).
Aplicar as relações métricas, incluindo as leis do seno e do cosseno ou as noções de congruência e semelhança, para resolver e elaborar \textbf{problemas} que envolvem triângulos, em variados contextos. ( EM13MAT309 ).
Resolver e elaborar \textbf{problemas} que envolvem o cálculo de áreas totais e de volumes de prismas, pirâmides e corpos redondos em situações reais ( como o cálculo do gasto de material para revestimento ou pinturas de objetos cujos formatos sejam composições dos sólidos estudados ), com ou sem apoio de \textbf{tecnologias} \textbf{digitais}. ( EM13MAT310 ).
Resolver e elaborar \textbf{problemas} de contagem envolvendo agrupamentos ordenados ou não de elementos, por meio dos princípios multiplicativo e aditivo, recorrendo a estratégias diversas, como o diagrama de árvore. ( EM13MAT311 ).
Identificar e descrever o espaço amostral de \textbf{eventos} aleatórios, realizando contagem das possibilidades, para resolver e elaborar \textbf{problemas} que envolvem o cálculo da probabilidade. ( EM13MAT312 ).
Resolver e elaborar \textbf{problemas} que envolvem o cálculo de probabilidade de \textbf{eventos} em experimentos aleatórios sucessivos. ( EM13MAT313 ).
Utilizar, quando necessário, a notação científica para expressar uma medida, compreendendo as noções de algarismos significativos e algarismos duvidosos, e reconhecendo que toda medida é inevitavelmente acompanhada de \textbf{erro}. ( EM13MAT314 ).
Resolver e elaborar \textbf{problemas} que envolvem grandezas determinadas pela \textbf{razão} ou pelo produto de outras ( velocidade, densidade demográfica, energia elétrica etc. ). ( EM13MAT315 ).
Investigar e registrar, por meio de um fluxograma, quando possível, um algoritmo que resolve um \textbf{problema}. ( EM13MAT316 ).
Resolver e elaborar \textbf{problemas}, em diferentes contextos, que envolvem cálculo e interpretação das medidas de tendência \textbf{central} ( média, moda, mediana ) e das medidas de dispersão ( amplitude, variância e desvio padrão ).
Objetos do conhecimento - Sistemas de equações lineares.
Gráficos de funções lineares com uma ou duas variáveis. - Função polinomial do 1º grau.
Função polinomial do 2º grau.
Variação entre grandezas ( proporcionalidade e não proporcionalidade ). - Conceitos de Matemática Financeira.
Juros \textbf{simples} e juros compostos.
Funções e gráficos de funções de 1º grau e exponencial. - Funções exponenciais.
Variação exponencial entre grandezas.
Noções de Matemática Financeira. - Logaritmo ( decimal e natural ).
Função logarítmica.
Variação entre grandezas : relação entre variação exponencial e logarítmica. - Trigonometria no triângulo retângulo ( principais \textbf{razões} trigonométricas ).
Trigonometria no ciclo trigonométrico.
Unidades de medidas de ângulos ( radianos ).
Funções trigonométricas ( função seno e função cosseno ). - Áreas de \textbf{figuras} geométricas ( cálculo por decomposição, composição ou aproximação ).
Expressões algébricas. - Lei dos senos e lei dos cossenos.
Congruência de triângulos ( por transformações geométricas – isometrias ).
Semelhança entre triângulos ( por transformações geométricas – homotetias ). - Geometria Métrica : poliedros e corpos redondos.
Área total e volume de prismas, pirâmides e corpos redondos. - Noções de combinatória : agrupamentos ordenáveis ( arranjos ) e não ordenáveis ( combinações ).
Princípio multiplicativo e princípio aditivo.
Modelos para contagem de dados : diagrama de árvore, listas, esquemas, desenhos etc. - Noções de probabilidade básica : espaço amostral, \textbf{evento} aleatório ( equiprovável ).
Contagem de possibilidades.
Cálculo de probabilidades \textbf{simples}. - \textbf{Eventos} dependentes e independentes.
Cálculo de probabilidade de \textbf{eventos} relativos a experimentos aleatórios sucessivos. - Notação científica.
Algarismos significativos e técnicas de arredondamento.
Estimativa e comparação de valores em notação científica e em arredondamentos.
Noção de \textbf{erro} em medições. - Grandezas determinadas pela \textbf{razão} ou produto de outras ( velocidade, densidade de um corpo, densidade demográfica, potência elétrica, bytes por segundo etc. ).
Conversão entre unidades compostas. - Noções básicas de Matemática Computacional.
Algoritmos e sua representação por fluxogramas. - Noções de estatística descritiva.
Medidas de tendência \textbf{central} : média, moda e mediana.
Medidas de dispersão : amplitude, variância e desvio-padrão. \textbf{Competência} Específica 4 Compreender e utilizar, com flexibilidade e precisão, diferentes registros de representação matemáticos ( algébrico, geométrico, estatístico, computacional etc. ), na busca de solução e comunicação de resultados de \textbf{problemas}.
A \textbf{competência} 4 complementa as demais no sentido de que utilizar, interpretar e resolver situações-problema realiza -se pela comunicação das ideias dos estudantes por meio da linguagem matemática.
Transitar entre os diversos tipos de representações ( simbólica, algébrica, gráfica, textual etc. ) permite a \textbf{compreensão} mais profunda dos conceitos e ideias da matemática.
Assim, a representação de uma mesma situação de diferentes formas estabelece conexões que possibilitam resolver \textbf{problemas} matemáticos, usando estratégias diversas.
Além disso, a \textbf{capacidade} de elaborar modelos matemáticos para expressar situações \textbf{implica} e revela a aprendizagem, além de potencializar o \textbf{letramento} matemático.
Essa \textbf{competência} está relacionada ao desenvolvimento das \textbf{Competências} Gerais 4 e 5 da BNCC, uma vez que a linguagem utilizada de modo flexível permite expressar ideias e informações que facilitam o entendimento e ampliar o repertório de formas de expressão, inclusive a \textbf{digital}, com espaço para \textbf{autoria} pessoal e criatividade do estudante.
Habilidades ( EM13MAT401 ).
Converter representações algébricas de funções polinomiais de 1º grau em representações geométricas no plano cartesiano, distinguindo os casos nos quais o comportamento é proporcional, recorrendo ou não a softwares ou \textbf{aplicativos} de álgebra e geometria dinâmica. ( EM13MAT402 ).
Converter representações algébricas de funções polinomiais de 2º grau em representações geométricas no plano cartesiano, distinguindo os casos nos quais uma variável for diretamente proporcional ao quadrado da outra, recorrendo ou não a softwares ou \textbf{aplicativos} de álgebra e geometria dinâmica, entre outros materiais. ( EM13MAT403 ).
Analisar e estabelecer relações, com ou sem apoio de \textbf{tecnologias} \textbf{digitais}, entre as representações de funções exponencial e logarítmica expressas em tabelas e em plano cartesiano, para identificar as características fundamentais ( domínio, imagem, crescimento ) de cada função. ( EM13MAT404 ).
Analisar funções definidas por uma ou mais sentenças ( tabela do Imposto de Renda, contas de luz, água, gás etc. ), em suas representações algébrica e gráfica, identificando domínios de validade, imagem, crescimento e decrescimento, e convertendo essas representações de uma para outra, com ou sem apoio de \textbf{tecnologias} \textbf{digitais}. ( EM13MAT405 ).
Utilizar conceitos iniciais de uma linguagem de programação na implementação de algoritmos escritos em linguagem corrente e/ou matemática. ( EM13MAT406 ).
Construir e interpretar tabelas e gráficos de frequências com base em dados obtidos em pesquisas por amostras estatísticas, incluindo ou não o uso de softwares que \textbf{inter-relacionem} estatística, geometria e álgebra. ( EM13MAT407 ).
Interpretar e comparar conjuntos de dados estatísticos por meio de diferentes diagramas e gráficos ( histograma, de caixa ( box-plot ), de ramos e folhas, entre outros ), reconhecendo os mais eficientes para sua \textbf{análise}.
Objetos do conhecimento - Funções afins, lineares, constantes.
Gráficos de funções a partir de transformações no plano.
Proporcionalidade : estudo do crescimento e variação de funções.
Estudo da variação de funções polinomiais de 1º grau : crescimento, decrescimento, taxa de variação da função. - Funções polinomiais de 2º grau.
Gráficos de funções a partir de transformações no plano.
Estudo do comportamento da função quadrática ( intervalos de crescimento/decrescimento, ponto de máximo/mínimo e variação da função ). - Funções : exponencial e logarítmica.
Gráfico de funções a partir de transformações no plano.
Estudo do crescimento e \textbf{análise} do comportamento das funções exponencial e logarítmica em intervalos numéricos. - Funções definidas por partes.
Gráficos de funções expressas por diversas sentenças. \textbf{Análise} do comportamento de funções em intervalos numéricos. - Noções elementares de matemática computacional : \textbf{sequências}, laços de repetição, variáveis e condicionais.
Algoritmos : modelagem de \textbf{problemas} e de soluções.
Linguagem da programação : fluxogramas. - Amostragem.
Gráficos e diagramas estatísticos : histogramas, polígonos de frequências.
Medidas de tendência \textbf{central} e medidas de dispersão. - Gráficos e diagramas estatísticos : histogramas, polígonos de frequências, diagrama de caixa, ramos e folhas etc.
Medidas de tendência \textbf{central} e medidas de dispersão. \textbf{Competência} Específica 5 Investigar e estabelecer conjecturas a respeito de diferentes conceitos e propriedades matemáticas, empregando estratégias e recursos, como observação de padrões, experimentações e diferentes \textbf{tecnologias}, identificando a necessidade, ou não, de uma demonstração cada vez mais formal na validação das referidas conjecturas.
A \textbf{competência} 5 tem como objetivo principal permitir que os estudantes se apropriem da forma de pensar matemática, como \textbf{ciência} com uma forma específica de validar suas conclusões pelo \textbf{raciocínio} lógico-dedutivo.
Não se trata de trazer para o Ensino Médio a Matemática formal dedutiva, mas de permitir que os jovens percebam a diferença entre uma dedução originária da observação empírica e uma dedução formal.
Além disso, é importante \textbf{verificar} que essa \textbf{competência} e suas habilidades não se desenvolvem em separado das demais ; ela é um \textbf{foco} a mais de atenção para o ensino em \textbf{termos} de formação dos estudantes, de modo que identifiquem a Matemática diferenciada das demais Ciências.
Diante disso, as habilidades para essa \textbf{competência} \textbf{demandam} que o estudante vivencie a investigação, a formulação de hipóteses e a tentativa de validação de suas hipóteses.
De \textbf{certa} forma, a proposta é que o estudante do Ensino Médio possa conhecer parte do processo de construção da Matemática, tal qual aconteceu ao longo da história, fruto do pensamento de muitos em diferentes culturas.
Um ponto de atenção está no fato de que algumas habilidades propostas pela BNCC ( 2018 ) para o desenvolvimento dessa \textbf{competência} \textbf{remetem} a conteúdos muito específicos, de pouca \textbf{aplicabilidade} e de difícil \textbf{contextualização}, no entanto favorecem a investigação e a formulação de hipóteses antes que os estudantes conheçam os conceitos ou a \textbf{teoria} subjacente a esses conteúdos específicos.
Tais habilidades possuem níveis diferentes de complexidade cognitiva, desde a identificação de uma propriedade até a investigação completa com dedução de uma regra ou procedimento.
Dessa forma, essa \textbf{competência} se relaciona com as \textbf{Competências} Gerais 2, 4, 5 e 7 da BNCC, uma vez que há o incentivo ao exercício da curiosidade intelectual na investigação, neste caso, com maior centralidade no conhecimento matemático.
A linguagem e os recursos \textbf{digitais} são ferramentas básicas e essenciais para facilitar a observação de regularidades, expressar ideias e construir \textbf{argumentos} com base em fatos.
Habilidades ( EM13MAT501 ).
Investigar relações entre números expressos em tabelas para representá -los no plano cartesiano, identificando padrões e criando conjecturas para generalizar e expressar algebricamente essa generalização, reconhecendo quando essa representação é de função polinomial de 1º grau. ( EM13MAT502 ).
Investigar relações entre números expressos em tabelas para representá -los no plano cartesiano, identificando padrões e criando conjecturas para generalizar e expressar algebricamente essa generalização, reconhecendo quando essa representação é de função polinomial de 2º grau do tipo y = ax². ( EM13MAT503 ).
Investigar pontos de máximo ou de mínimo de funções quadráticas em contextos envolvendo superfícies, Matemática Financeira ou Cinemática, entre outros, com apoio de \textbf{tecnologias} \textbf{digitais}. ( EM13MAT504 ).
Investigar processos de obtenção da medida do volume de prismas, pirâmides, cilindros e cones, incluindo o princípio de Cavalieri, para a obtenção das fórmulas de cálculo da medida do volume dessas \textbf{figuras}. ( EM13MAT505 ).
Resolver \textbf{problemas} sobre ladrilhamento do plano, com ou sem apoio de \textbf{aplicativos} de geometria dinâmica, para conjecturar a respeito dos tipos ou composição de polígonos que podem ser utilizados em ladrilhamento, generalizando padrões observados. ( EM13MAT506 ).
Representar graficamente a variação da área e do perímetro de um polígono regular quando os comprimentos de seus \textbf{lados} variam, analisando e classificando as funções envolvidas. ( EM13MAT507 ).
Identificar e associar \textbf{progressões} aritméticas ( PA ) a funções afins de domínios discretos, para \textbf{análise} de propriedades, dedução de algumas fórmulas e resolução de \textbf{problemas}. ( EM13MAT508 ).
Identificar e associar \textbf{progressões} geométricas ( PG ) a funções exponenciais de domínios discretos, para \textbf{análise} de propriedades, dedução de algumas fórmulas e resolução de \textbf{problemas}. ( EM13MAT509 ).
Investigar a deformação de ângulos e áreas provocada pelas diferentes \textbf{projeções} usadas em cartografia ( como a cilíndrica e a cônica ), com ou sem suporte de \textbf{tecnologia} \textbf{digital}. ( EM13MAT510 ).
Investigar conjuntos de dados relativos ao comportamento de duas variáveis numéricas, usando ou não \textbf{tecnologias} da informação, e, quando apropriado, levar em conta a variação e utilizar uma reta para descrever a relação observada. ( EM13MAT511 ).
Reconhecer a existência de diferentes tipos de espaços amostrais, discretos ou não, e de \textbf{eventos}, equiprováveis ou não, e investigar implicações no cálculo de probabilidades.
Objetos do conhecimento - Funções polinomiais do 1º grau ( função afim, função linear, função constante, função identidade ).
Gráficos de funções.
Taxa de variação de funções polinomiais do 1º grau. - Funções polinomiais do 2º grau ( função quadrática ) : gráfico, raízes, pontos de máximo/mínimo, crescimento/decrescimento, concavidade.
Gráficos de funções.
Funções polinomiais do 2º grau ( função quadrática ).
Gráficos de funções.
Pontos \textbf{críticos} de uma função quadrática : concavidade, pontos de máximo ou de mínimo.
Sólidos geométricos ( prismas, pirâmides, cilindros e cones ).
Cálculo de volume de sólidos geométricos. - Polígonos regulares e suas características : ângulos internos, ângulos externos etc.
Pavimentações no plano ( usando o mesmo tipo de polígono ou não ).
Linguagem algébrica : fórmulas e habilidade de generalização. - Polígonos regulares ( perímetro e área ).
Funções ( linear e quadrática ).
Funções afins. - \textbf{Sequências} numéricas : \textbf{progressões} aritméticas ( P.A. ).
Função exponencial. \textbf{Sequências} numéricas : \textbf{progressões} geométricas ( P.G. ).
Transformações geométricas ( isometrias e homotetias ).
Posição de \textbf{figuras} geométricas ( tangente, secante, externa ).
Inscrição e circunscrição de sólidos geométricos.
Noções básicas de cartografia ( \textbf{projeção} cilíndrica e cônica ). - Funções polinomiais do 1º grau ( função afim, linear e constante ).
Gráficos de funções.
Taxa de variação de uma função ( crescimento/decrescimento ). \textbf{Razões} trigonométricas : tangente de um ângulo.
Equação da reta : coeficiente angular. - Probabilidade.
Espaços amostrais discretos ou contínuos. \textbf{Eventos} equiprováveis ou não equiprováveis.
Em continuidade à apresentação dos conceitos e \textbf{aplicabilidade} das \textbf{competências} e das habilidades, vale \textbf{salientar} a importância da \textbf{progressão} das habilidades com o \textbf{intuito} de \textbf{auxiliar} os estudantes na aquisição dos objetos de conhecimento de forma gradativa, contemplando os objetivos de aprendizagem dos mais \textbf{simples} até os mais \textbf{complexos}.
Assim, para contribuir com o planejamento e \textbf{compreensão} dos docentes, as habilidades da área de Matemática e suas Tecnologias foram distribuídas de \textbf{maneira} a propiciar condições de aprendizagem a fim de ampliar, aprofundar e consolidar as unidades curriculares da área. \textbf{Posto} isso, a distribuição das habilidades nos três anos apresentadas no Quadro Organizador da Área de Matemática e suas Tecnologias para o Ensino Médio considerou a Taxonomia de Bloom com a finalidade de contribuir com a aquisição dos objetos de conhecimento mediante à complexidade de cada habilidade.
Dessa forma, ressalta -se que há possibilidades de retomada das habilidades nos outros anos, para assegurar que o estudante tenha a oportunidade de desenvolver habilidades e \textbf{competências} não contempladas em sua integralidade no ano anterior.
Logo, todas as habilidades podem ser \textbf{mobilizadas} por todos os anos do Ensino Médio, sendo oportunizado ao docente o reposicionamento.
No entanto, é \textbf{imprescindível} que no \textbf{decorrer} de cada ano sejam contempladas as habilidades conforme o \textbf{quadro} 1 apresentado abaixo.
No Quadro 1, buscou -se equilibrar os objetivos de aprendizagem ao realizar a distribuição das habilidades específicas da área de Matemática e suas Tecnologias a fim de contribuir para a formação integral dos estudantes, pois perpassarão por todas as habilidades no \textbf{decorrer} do Ensino Médio.
Vale \textbf{salientar} que nos três anos, o estudo das cinco \textbf{competências} específicas perpassa de \textbf{maneira} progressiva, contemplando todas as habilidades que envolvem as unidades curriculares.
Assim, em continuidade aos estudos dos conceitos relacionados ao Ensino Médio e suas \textbf{aplicabilidades} serão demonstrados quais são os objetivos e as expectativas de aprendizagem para a área de Matemática e suas Tecnologias.

% matched lemmas: análise, aplicabilidade, aplicativo, argumento, autoria, auxiliar, capacidade, central, certo, ciência, competência, complexo, compreensão, contemporâneo, contextualização, crítico, decorrer, demanda, demandar, digital, divulgar, econômico, erro, evento, exigir, figura, foco, globalizado, implicar, imprescindível, induzir, infográfico, inter-relacionar, intuito, lado, letramento, maneira, mero, mobilizar, método, postar, problema, progressão, projeção, quadro, raciocínio, razão, relevante, remeter, restringir, salientar, sequência, simples, tanto, tecnologia, teoria, termos, verificar
\end{textsample}
